\chapter{BẤT PHƯƠNG TRÌNH VÀ HỆ BẤT PHƯƠNG TRÌNH BẬC NHẤT HAI ẨN}
\section{Bất phương trình bậc nhất hai ẩn}

\subsection{LÝ THUYẾT TRỌNG TÂM}

\begin{kttrongtam}
	\textbf{\textit{Bất phương trình bậc nhất hai ẩn}} $x, y$ là bất phương trình có một trong các dạng
	$$ax+by+c<0,\; ax+by+c>0 ,\; ax+by+c \leq 0 ,\; a x+b y+c \geq 0$$
	trong đó $a$, $b$, $c$ là những số cho trước; $a$, $b$ không đồng thời bằng $0$ và $x$, $y$ là các ẩn.
\end{kttrongtam}

%===========VD1
\begin{vd}%[0D2N1-1]
	Tìm bất phương trình bậc nhất hai ẩn trong các bất phương trình sau đây
	\begin{enumEX}{2}
		\item[a)] $x-5y+2 < 0$;
		\item[b)] $9x^2+8y-7 \geq 0$;
		\item[c)] $3x-2>0$;
		\item[d)] $4y+11 \leq 0$.
	\end{enumEX}
	\loigiai{
		Các bất phương trình $x-5y+2 < 0$, $3x-2>0$, $4y+11 \leq 0$ là các bất phương trình bậc nhất hai ẩn.\\
		Bất phương trình $9x^2+8y-7 \geq 0$ không là bất phương trình bậc nhất hai ẩn vì có chứa $x^2$.
	}
\end{vd}
%==========VD2
\begin{vd}%[0D2N1-1]
 	Tìm các bất phương trình bậc nhất hai ẩn trong các bất phương trình sau đây
 	\begin{enumEX}{2}
		\item[a)] $\dfrac{-1}{7}x-\dfrac{y^2}{3} \le 8$;
		\item[b)] $\sqrt{2} x^2-5 \sqrt{y} \ge 8$;
		\item[c)] $2 \cdot \dfrac{1}{x}-5\cdot \dfrac{1}{y}>8$;
		\item[d)] $-\dfrac{2}{5} x-5^2 y \leq-\sqrt{15}$.
	\end{enumEX}
 	\loigiai{
 	Bất phương trình bậc nhất hai ẩn là  $-\dfrac{2}{5} x-5^2 y \leq-\sqrt{15}$.
 	}
 \end{vd}
\begin{kttrongtam}
	Xét bất phương trình $ax+by+c<0$.\\
	Mỗi cặp số $(x_0;y_0)$ thỏa mãn $ax_0+by_0+c<0$ được gọi là một \textbf{\textit{nghiệm}} của bất phương trình đã cho.
\end{kttrongtam}

\begin{note}
	Nghiệm của bất phương trình $ax+by+c>0$, $ax+by+c \leq 0$, $ax+by+c \geq 0$ được định nghĩa tương tự.
\end{note}

%==========VD3
\begin{vd}%[0D2N1-2]
	Cặp số nào sau đây là nghiệm của bất phương trình $20x+50y-700 < 0$?
	\begin{enumEX}{2}
		\item[a)] $(5;6)$;
		\item[b)] $(9;11)$.
	\end{enumEX}
	\loigiai{
		\begin{enumerate}
			\item[a)] Vì $20 \cdot 5 + 50 \cdot 6 - 700=-300<0$ nên $(5;6)$ là một nghiệm của bất phương trình $20x+50y-700 < 0$.
			\item[b)] Vì $20 \cdot 9 +50 \cdot 11 - 700=30 >0$ nên $(9;11)$ không phải là một nghiệm của bất phương trình $20x+50y-700 < 0$.
		\end{enumerate}
	}
\end{vd}
%==========VD4
\begin{vd}%[0D2H1-2]
	Điểm $O(0 ; 0)$ thuộc miền nghiệm của bất phương trình nào dưới đây?
	\begin{enumEX}{2}
		\item[a)] $x+3 y+2 \le 0$;
		\item[b)] $x+y+2 \le 0$;
		\item[c)] $2 x+5 y-2<0$;
		\item[d)] $2 x+y<0$.
	\end{enumEX}
	\loigiai{
	\begin{itemize}
		\item Xét bất phương trình $x+3 y+2 \le 0$, ta có $0 + 3\cdot 0 + 2 = 2>0$. Do đó điểm $O(0 ; 0)$ không  thuộc miền nghiệm của bất phương trình $x+3 y+2 \le 0$.
		\item Xét bất phương trình $x+y+2 \le 0$, ta có $0 + 0 + 2 = 2>0$. Do đó điểm $O(0 ; 0)$ không  thuộc miền nghiệm của bất phương trình $x+3 y+2 \le 0$.
		\item Xét bất phương trình $2 x+5 y-2<0$, ta có $2\cdot0 + 5\cdot 0 -2=-2 < 0$. Do đó điểm $O(0 ; 0)$  thuộc miền nghiệm của bất phương trình $2 x+5 y-2<0$.
		\item Xét bất phương trình $2 x+y<0$, ta có $2\cdot0 +  0  = 0$. Do đó điểm $O(0 ; 0)$ không  thuộc miền nghiệm của bất phương trình $2 x+y<0$.
	\end{itemize}	
	Vậy điểm $O(0 ; 0)$  thuộc miền nghiệm của bất phương trình $2 x+5 y-2<0$.
	}
\end{vd}
%=============VD5
\vspace{0.5cm}
\begin{kttrongtam}
	Trong mặt phẳng tọa độ $Oxy$, tập hợp các điểm $(x_0;y_0)$ sao cho $ax_0+by_0+c<0$ được gọi là \textbf{\textit{miền nghiệm}} của bất phương trình $ax+by+c<0$.
\end{kttrongtam}
Ta có thể biểu diễn miền nghiệm của bất phương trình bậc nhất hai ẩn $ax+by+c<0$ như sau:
\begin{itemize}
	\item \textbf{Bước 1:} Trên mặt phẳng $Oxy$, vẽ đường thẳng $\Delta\colon ax+by+c=0$.
	\item \textbf{Bước 2:} Lấy một điểm $(x_0;y_0)$ không thuộc $\Delta$. Tính $ax_0+by_0+c$.
	\item \textbf{Bước 3:} Kết luận
	\begin{itemize}
		\item Nếu $ax_0+by_0+c<0$ thì miền nghiệm của bất phương trình đã cho là nửa mặt phẳng (không kể bờ $\Delta$) chứa điểm $(x_0;y_0)$.
		\item Nếu $ax_0+by_0+c>0$ thì miền nghiệm của bất phương trình đã cho là nửa mặt phẳng (không kể bờ $\Delta$) không chứa điểm $(x_0;y_0)$.
	\end{itemize}
\end{itemize}

\begin{note}
	Đối với các bất phương trình bậc nhất hai ẩn dạng $ax+by+c \le 0$ (hoặc $ax+by+c \ge 0$) thì miền nghiệm là miền nghiệm của bất phương trình $ax+by+c < 0$ (hoặc $ax+by+c > 0$) kể cả bờ.
\end{note}

%==========VD5
\begin{vd}%[0D2N1-2]
	Biểu diễn miền nghiệm của các bất phương trình sau:
	\begin{enumEX}{2}
		\item[a)] $x-2y-1>0$;
		\item[b)] $x+y-1 \le 0$.
	\end{enumEX}
	\loigiai{
		\begin{enumerate}
			\item[a)] $x-2y-1>0$.
			\immini{
				Vẽ đường thẳng $\Delta \colon x-2y-1=0$ đi qua hai điểm $A(1;0)$ và $B\left(0;-\dfrac{1}{2}\right)$.\\
				Xét gốc tọa độ $O(0;0)$.\\
				Ta thấy $O \notin \Delta$ và $0-2 \cdot 0 -1<0$.\\
				Do đó miền nghiệm của bất phương trình là nửa mặt phẳng không kể bờ $\Delta$, không chứa gốc tọa độ $O$ (miền không gạch chéo trên hình bên).
			}{\vspace{-0.5cm}
				\begin{tikzpicture}[scale=1, font=\footnotesize, line join=round, line cap=round, >=stealth]
					\draw[->] (-2,0)--(3.8,0) node [below]{$x$};
					\draw[->] (0,-2)--(0,1.8) node [left]{$y$};
					\node at (0,0) [above left=-3pt]{$O$};
					\clip (-2,-2) rectangle (3.8,1.8);
					\fill[pattern=north west lines,opacity=0.6] plot[domain=-3.6:3.8](\x,{1/2*(\x)-1/2})--(3.8,1.8)--(-3.6,1.8)--cycle;
					\draw plot[domain=-3.6:3.8](\x,{1/2*(\x)-1/2})node[shift={(-120:10pt)}]{$\Delta$};
					\draw[fill=black] (1,0) circle(1pt) node[shift={(90:10pt)}]{$A$};
					\draw[fill=black] (0,-1/2)circle(1pt) +(-35:8pt)node{$B$};
					\node at (1,0)[shift={(-90:8pt)}]{$1$};
					\node at (0,-1/2)[shift={(-125:13pt)}]{$-\tfrac{1}{2}$};
				\end{tikzpicture}
			}
			\item[b)] $x+y-1\le 0$. 
			\immini{
				Vẽ đường thẳng $\Delta \colon x+y-1=0$ đi qua hai điểm $A(1;0)$ và $B(0;1)$.\\
				Xét gốc tọa độ $O(0;0)$.\\
				Ta thấy $O \notin \Delta $ và $0+0-1<0$.\\
				Do đó miền nghiệm của bất phương trình là nửa mặt phẳng kể cả bờ $\Delta$, chứa gốc tọa độ $O$ (miền không gạch chéo trên hình bên).
			}{\vspace{-0.5cm}
				\begin{tikzpicture}[scale=1, font=\footnotesize, line join=round, line cap=round, >=stealth]
					\draw[->] (-2,0)--(3,0) node [below]{$x$};
					\draw[->] (0,-2)--(0,1.8) node [left]{$y$};
					\node at (0,0) [below left=-3pt]{$O$};
					\clip (-2,-2) rectangle (3,1.8);
					\fill[pattern=north east lines,opacity=0.6] plot[domain=-3.6:3.8](\x,{-(\x)+1})--(3.8,1.8)--cycle;
					\draw plot[domain=-3.6:3](\x,{-(\x)+1})node[shift={(153:15pt)}]{$\Delta$};
					\draw[fill=black] (1,0)circle(1pt) +(65:8pt)node{$A$};
					\draw[fill=black] (0,1)circle(1pt) +(15:8pt)node{$B$};
					\node at (1,0)[shift={(-90:8pt)}]{$1$};
					\node at (0,1)[shift={(185:8pt)}]{$1$};
				\end{tikzpicture}
			}
		\end{enumerate}
	}
\end{vd}
%=========VD6
\begin{vd}%[0D2N1-2]
	Biểu diễn miền nghiệm của bất phương trình $2x-3y \leq 0$ trên mặt phẳng toạ độ.
	\loigiai{
		\immini{
			\begin{itemize}
				\item \textbf{Bước 1.} Vẽ đường thẳng $d\colon 2x-3y=0$ trên mặt phẳng toạ độ $Oxy$.
				\item \textbf{Bước 2.} Lấy điểm $M_{0} (0;1)$ không thuộc $d$ và thay $x=0, y=1$ vào biểu thức $2x-3y$ ta được $2 \cdot 0 - 3 \cdot 1 =-3<0$.
			\end{itemize}
			Do đó miền nghiệm của bất phương trình là nửa mặt phẳng bờ $d$ chứa điểm $M_{0}$ (miền không bị gạch).
		}{
			\begin{tikzpicture}[scale=0.7,line join=round, line cap=round,>=stealth,thick]
				\tikzset{label style/.style={font=\footnotesize}}
				\begin{scope}
					\clip (-1,-1) rectangle (6,4.3);
					\fill[pattern=north west lines,opacity=0.4] (-3,-2)--(6,-2)--(6,4)--cycle;
					\draw[dashed] (0,2)--(3,2)--(3,0);
					\draw [thick](-3,-2)--(6,4) node [pos=0.8, above, sloped] {$d\colon 2x-3y=0$};
				\end{scope}
				\draw[->] (-1,0)--(6.1,0) node[below]{$x$};
				\draw[->] (0,-1)--(0,4) node[left]{$y$};
				\draw (0,0) node[below right]{$O$};
				\foreach \x in {1,2,3}
				\draw[thin] (\x,1pt)--(\x,-1pt) node [below] {$\x$};
				\foreach \y in {1,2}
				\draw[thin] (1pt,\y)--(-1pt,\y) node [left] {$\y$};
				\fill[red] (0,1) circle (1pt) node [right] {$M_{0}$};
			\end{tikzpicture}
		}
	}
\end{vd}


\subsection{BÀI TẬP RÈN LUYỆN}

\caulc
%========1
\begin{ex}%[0D2N1-1]
 	Bất phương trình nào sau đây là bất phương trình bậc nhất hai ẩn?
 	\choice
 	{$\dfrac{-1}{7}x-\dfrac{y^2}{3} \le 8$ }
 	{$\sqrt{2} x^2-5 \sqrt{y} \ge 8$}
 	{$2 \cdot \dfrac{1}{x}-5\cdot \dfrac{1}{y}>8$}
 	{\True $-\dfrac{2}{5} x-5^2 y \leq-\sqrt{15}$}
 	\loigiai{
 		Bất phương trình bậc nhất hai ẩn là  $-\dfrac{2}{5} x-5^2 y \leq-\sqrt{15}$.
 	}
 \end{ex}
%========2
\begin{ex}%[0D2N1-1]
	Bất phương trình $3x-2(y-x+1)>0$ tương đương với bất phương trình nào sau đây?
	\choice
	{$x-2y-2>0$}
	{\True $5x-2y-2>0$}
	{$5x-2y-1>0$}
	{$4x-2y-2>0$}
	\loigiai{
		Ta có $3x-2(y-x+1)>0\Leftrightarrow 3x-2y+2x-2>0\Leftrightarrow 5x-2y-2>0$.
	}
\end{ex}
%========3
\begin{ex}%[0D2N1-1]
 	Điểm $O(0 ; 0)$ thuộc miền nghiệm của bất phương trình nào dưới đây?
 	\choice
 	{$x+3 y+2 \le 0$}
 	{$x+y+2 \le 0$}
 	{\True $2 x+5 y-2<0$}
 	{$2 x+y<0$}
 	\loigiai{
 		Thay tọa độ $O(0;0)$ vào các bất phương trình:
 		\begin{itemize}
 			\item[•] Với bất phương trình $x+3 y+2 \le 0$, ta có $0+3\cdot 0+2\le 0$ sai.
 			\item[•] Với bất phương trình $x+y+2 \le 0$, ta có $0+0+2\le 0$ sai.
 			\item[•] Với bất phương trình $2 x+5 y-2<0$, ta có $2\cdot 0+5\cdot 0-2<0$ đúng.
 			\item[•] Với bất phương trình $2 x+y<0$, ta có $2\cdot 0+0>0$ sai.
 		\end{itemize}
 		Vậy $O(0;0)$ thuộc miền nghiệm bất phương trình $2 x+5 y-2<0$.
 	}
 \end{ex}
%========4
\begin{ex}%[0D2N1-1]
	Điểm $A(-1;3)$ thuộc miền nghiệm của bất phương trình
	\choice
	{$x+3y<0$}
	{$3x-y>0$}
	{\True  $-3x+2y-4>0$}
	{$2x-y+4>0$}
	\loigiai{
		Thay tọa độ $A(-1;3)$ vào các bất phương trình:
		\begin{itemize}
			\item[•] Với bất phương trình $x+3y<0$, ta có $(-1)+3\cdot 3<0$ sai.
			\item[•] Với bất phương trình $3x-y>0$, ta có $3\cdot (-1)-3>0$ sai.
			\item[•] Với bất phương trình $-3x+2y-4>0$, ta có $-3\cdot (-1)+2\cdot 3-4>0$ đúng.
			\item[•] Với bất phương trình $2x-y+4>0$, ta có $2\cdot (-1)-3+4>0$ sai.
		\end{itemize}
		Vậy $A(-1;3)$ thuộc miền nghiệm bất phương trình $-3x+2y-4>0$.
	}
\end{ex}
%========5
\begin{ex}%[0D2N1-1]
	Trong các bất phương trình sau đây, đâu là bất phương trình bậc nhất hai ẩn
	\choice
	{$2x^2-3x \geq 1$}
	{\True $2x+y\leq 1$}
	{$3x+\sqrt{y} \leq 0$}
	{$3x+y=1$}
	\loigiai{
	Theo định nghĩa $2x+y\leq 1$ là bất phương trình bậc nhất hai ẩn.
	}
\end{ex}
%========6
\begin{ex}%[0D2N1-2]
	Trong các bất phương trình sau, bất phương trình nào là bất phương trình bậc nhất hai ẩn?
	\choice
	{$2x^2+3y>4$}
	{$xy+x<5$}
	{\True $3^2x+4^3y \geq 6$}
	{$x+y^3 \leq 3$}
	\loigiai{
	Bất phương trình $3^2x+4^3y \geq 6$ là bất phương trình bậc nhất hai ẩn.
	}
\end{ex}
%========12
\begin{ex}%[0D2N1-1]
	Cho bất phương trình $2x+3y-6 \leq 0\quad (1)$. Chọn khẳng định đúng trong các khẳng định sau
	\choice
	{Bất phương trình $(1)$ chỉ có một nghiệm duy nhất}
	{Bất phương trình $(1)$ vô nghiệm}
	{\True Bất phương trình $(1)$ có vô số nghiệm}
	{Bất phương trình $(1)$ có tập nghiệm là $\mathbb{R}$}
	\loigiai{
	Trên mặt phẳng tọa độ, đường thẳng $(d) \colon 2x+3y-6=0$ chia mặt phẳng thành hai nửa mặt phẳng.\\
	Chọn điểm $O(0;0)$ không thuộc đường thẳng đó. Ta thấy $(x;y)=(0;0)$ là nghiệm của bất phương trình đã cho.\\ Vậy miền nghiệm của bất phương trình là nửa mặt phẳng bờ $(d)$ chứa điểm $O(0;0)$ kể cả $(d)$.\\
	Vậy bất phương trình $(1)$ có vô số nghiệm.}
\end{ex}
%========7
\begin{ex}%[0D2N1-2]
 	Cho bất phương trình bậc nhất hai ẩn $2 x-5 y \leq 8$. Khẳng định nào dưới đây là khẳng định \textbf{sai}?
 	\choice
 	{$(3 ;-4)$ là một nghiệm của bất phương trình}
 	{\True $(-2 ; 2)$ không là một nghiệm của bất phương trình}
 	{ $(-3 ;-1)$ là một nghiệm của bất phương trình}
 	{ $(5 ; 0)$ không là một nghiệm của bất phương trình}
 	\loigiai{
 		Thay $(-2 ; 2)$ vào bất phương trình $2 x-5 y \leq 8$ ta được $2 \cdot (-2)-5\cdot 2 \leq 8 \Leftrightarrow-14 \leq 8$
 		(đúng), nên $(-2 ; 2)$ là một nghiệm của bất phương trình $2 x-5 y \leq 8$.
 	}
 \end{ex}
%========8
\begin{ex}%[0D2N1-2]
 	Cho bất phương trình $-2 x+3 y>3$. Khẳng định nào dưới đây \textbf{sai}?
 	\choice 
 	{$(0 ; 0)$ không là nghiệm bất phương trình}
 	{\True $(-1 ; 1)$ không là nghiệm bất phương trình}
 	{$(0 ; 1)$ không là nghiệm bất phương trình}
 	{ $(1 ; 3)$ là nghiệm bất phương trình}
 	\loigiai{
 		Thay $x=-1$, $y=1$ vào bất phương trình $2+3>3$ (đúng).\\
 		Vì vậy $(-1 ; 1)$ là nghiệm bất phương trình.
 	}
 \end{ex}
%========9
\begin{ex}%[0D2N1-2]
 	Miền nghiệm của bất phương trình $3(x-1)+4(y-2)<5x-3$ là nửa mặt phẳng chứa điểm
 	\choice
 	{$Q(-5;3)$}
 	{\True $O(0;0)$}
 	{$N(-4;2)$}
 	{$P(-2;2)$}
 	\loigiai{
 		Ta có $3(x-1)+4(y-2)<5x-3\Leftrightarrow -x+2y-4<0 \quad(\ast)$.\\
 		Thay $O(0;0)$ vào $(\ast)$ ta được $-0+1\cdot (-4)= -4 < 0$.\\
 		Do đó $O(0;0)$ thuộc miền nghiệm của bất phương trình.
 	}
 \end{ex}
%========10
\begin{ex}%[0D2N1-2]
 	Điểm $O(0 ; 0)$ thuộc miền nghiệm của bất phương trình nào dưới đây?
 	\choice
 	{$3x+3y \le -2$}
 	{$\dfrac{1}{3}x+2y+2 \le 0$}
 	{\True $2x+4y-2<0$}
 	{$x+10y>0$}
 	\loigiai{
 		Thay $x=y=0$ vào bất phương trình $2x+4y-2<0$, ta được $2\cdot 0+4 \cdot 0 -2<0$ (đúng).\\
 		Vì vậy $O(0 ; 0)$ thuộc miền nghiệm.
 	}
 \end{ex}
%========12
\begin{ex}%[0D2H1-2]
	\immini{Miền không gạch chéo (không kể bờ $d$) trong hình bên là miền nghiệm của bất phương trình nào trong các bất phương trình dưới đây?
	\choice
	{$2x+3y<6$}
	{\True $2x+3y>6$}
	{$\dfrac{x}{2}+\dfrac{y}{3}>0$}
	{$\dfrac{x}{2}+\dfrac{y}{3}<1$}
	}
	{
	\begin{tikzpicture}[scale=0.8, font=\footnotesize, line join=round, line cap=round,>=stealth]
	\def \xmin{-1};
	\def \xmax{4};
	\def \ymin{-1};
	\def \ymax{3};
	\draw[->] (\xmin, 0.) -- (\xmax,0.) node[anchor=north] {$x$};
	\draw[->] (0.,\ymin) -- (0.,\ymax) node[anchor=west] {$y$};
	\clip(\xmin-0.1,\ymin-0.1) rectangle (\xmax+0.1,\ymax+0.1);
	\fill[pattern=north east lines,opacity=0.4] (-1,3) plot[smooth,domain=-1:4] (\x,{(6-2*(\x))/3})--(4,-1)--(-1,-1)--cycle;
	%\fill[pattern=north west lines,opacity=0.4] (0,0) plot[smooth,domain=0:1] (\x,{(\x)*(\x-1)})--(1,0)--cycle;
	\draw[smooth] plot[domain=\xmin:\xmax] (\x,{(6-2*(\x))/3});
	\draw[fill=black] (0,0) circle (1pt) node[below left] {$O$} (1,0) circle(1pt) node[below]{$1$} (2,0) circle(1pt) node[below]{$2$} (3,0) circle(1pt) node[below]{$3$} (0,2) circle(1pt) node[below left]{$2$} (0,1) circle(1pt) node[left]{$1$};
	\draw (1.5,1) node[right]{$d$};
	\end{tikzpicture}
	}
	\loigiai{
	Xét đường thẳng $d\colon y=ax+b$ đi qua hai điểm $A(3,0)$ và $B(0,2)$.\\
	Ta có $\heva{& 0=3a+b\\ & 2=b}\Leftrightarrow \heva{& a=-\dfrac{2}{3}\\& b=2}\Rightarrow d\colon y=-\dfrac{2}{3}x+2\Leftrightarrow d\colon 2x+3y-6=0$.\\
	Xét điểm $O(0;0)$, ta có $2\cdot 0+3\cdot 0-6=-6<0$ nên điểm $O$ thuộc miền nghiệm của bất phương trình $2x+3y-6<0$.\\
	Vậy miền không gạch chéo (không kể bờ $d$) là miền nghiệm của bất phương trình $2x+3y-6>0$ hay $2x+3y>6$.
	}
\end{ex}
%==============================================
\cauds
%========1
\begin{ex}%[0D2H1-2]
 	Cho điểm $(-1 ; 2)$ và các bất phương trình $3 x-5 y<-15$, $2 x+y \le 0$, $ 3 x-9 y>7$, $-4 x+3 y \ge 5$.
 	\choiceTF
 	{\True $(-1 ; 2)$ không là một nghiệm của bất phương trình $3 x-5 y<-15$}
 	{\True $(-1 ; 2)$ là một nghiệm của bất phương trình $2 x+y \le 0$}
 	{$(-1 ; 2)$ là một nghiệm của bất phương trình $3 x-9 y>7$}
 	{\True $(-1 ; 2)$ là một nghiệm của bất phương trình $-4 x+3 y \ge 5$}
 	\loigiai{
 		\begin{itemchoice}
 			\itemch
 			Thay $(-1 ; 2)$ vào bất phương trình $3 x-5 y<-15$ ta được \\
 			$3 \cdot ( -1)-5 \cdot 2<-15 \Leftrightarrow-13<-15$ (vô lí), nên $(-1 ; 2)$ không là một nghiệm của bất phương trình $3 x-5 y<-15$.
 			\itemch
 			Thay $(-1 ; 2)$ vào bất phương trình $2 x+y \leq 0$ ta được $2 \cdot(-1)+2 \leq 0 \Leftrightarrow 0 \leq 0$
 			(đúng), nên $(-1 ; 2)$ là một nghiệm của bất phương trình $2 x+y \leq 0$.
 			\itemch
 			Thay $(-1 ; 2)$ vào bất phương trình $3 x-9 y>7$ ta được $3 \cdot (-1)-9 \cdot 2>7 \Leftrightarrow-21>7$
 			(vô lí), nên $(-1 ; 2)$ không là một nghiệm của bất phương trình $3 x-9 y>7$.
 			\itemch
 			Thay $(-1 ; 2)$ vào bất phương trình $-4 x+3 y \geq 5$ ta được $-4 \cdot (-1)+3\cdot 2 \geq 5 \Leftrightarrow
 			10 \geq 5$ (đúng), nên $(-1 ; 2)$ là một nghiệm của bất phương trình $-4 x+3 y \geq 5$.
 	\end{itemchoice}}
 \end{ex}
%========2
\begin{ex}%[0D2H1-1]
 	Cho bất phương trình bậc nhất hai ẩn $x-2y+2 \leq 0$.
 	\choiceTF
 	{\True $(1;4)$ là nghiệm của bất phương trình $x-2y+2 \leq 0$}
 	{$(0;3)$ không là nghiệm của bất phương trình $x-2y+2 \leq 0$}
 	{Đường thẳng $d \colon x-2y+2=0$ đi qua hai điểm $A(0;-1)$ và $B(2;0)$}
 	{\True Miền nghiệm của bất phương trình $x-2y+2 \leq 0$ là nửa mặt phẳng kể cả bờ $x-2y+2=0$, không chứa điểm $(0;0)$}
 	\loigiai{
 		\begin{itemchoice}
 			\itemch Ta có $1 \cdot 1 - 2 \cdot 4 + 2 = -5<0$ nên $(1;4)$ là nghiệm của bất phương trình $x-2y+2 \leq 0$.
 			\itemch Ta có $1 \cdot 0 - 2 \cdot 3 + 2 = -3<0$ nên $(0;3)$ là nghiệm của bất phương trình $x-2y+2 \leq 0$.
 			\itemch Ta có $2-2\cdot 0+2 \neq 0$ nên điểm $B(2;0)$ không nằm trên đường thẳng $d$.
 			\itemch Xét đường thẳng $d \colon x-2y+2=0$ đi qua hai điểm $A(0;1)$ và $B(-2;0)$.\\
 			Xét gốc tọa độ $O(0;0) \notin d$ và ta có $1 \cdot 0 - 2 \cdot 0 + 2 = 2 >0$. \\
 			Do đó, miền  nghiệm của bất phương trình $x-2y+2 \leq 0$ là nửa mặt phẳng kể cả bờ $d$ không chứa gốc tọa độ $O$.
 		\end{itemchoice}
 	}
 \end{ex}
%========3
\begin{ex}%[0D2H1-3]
 	An thích ăn hai loại trái cây là cam và xoài, mỗi tuần mẹ cho An $200\,000$ đồng để mua trái cây. Biết rằng giá cam là $15\,000$ đồng/kg, giá xoài là $30\,000$ đồng/kg. Gọi $x$, $y$ lần lượt là số kilôgam cam và xoài mà An có thể mua về sử dụng trong một tuần.
 	\choiceTF 
 	{\True Trong tuần, số tiền An có thể mua cam là $15\,000x$, số tiền An có thể mua xoài là $30\,000 y$, ($x$, $y>0$)}
 	{Bất phương trình bậc nhất cho hai ẩn $x$, $y$ là $3 x+6 y \geq 40$}
 	{\True Cặp số $(5 ; 4)$ thỏa mãn bất phương trình bậc nhất cho hai ẩn $x$, $ y$}
 	{An có thể mua $4 $ kg cam, $5 $ kg xoài trong tuần}
 	\loigiai{
 		\begin{itemchoice}
 			\itemch Trong tuần, số tiền An có thể mua cam là $15\,000 x$, số tiền An có thể mua xoài là $30\,000 y$ ($x$, $y>0$).
 			\itemch Ta có bất phương trình $15\,000 x+30\,000 y \leq 200\,000 \Leftrightarrow 3 x+6 y \leq 40$.\quad $(\ast)$
 			\itemch Xét $x=5$, $y=4$, thay vào bất phương trình $3\cdot5+6\cdot 4 \leq 40$ (đúng) nên $(5 ; 4)$ là một nghiệm của $(\ast)$.
 			\itemch Xét $x=4$, $y=5$, thay vào bất phương trình $3\cdot4+6\cdot 5 \leq 40$ (sai) nên $(4 ; 5)$ không là một nghiệm của $(\ast)$.
 		\end{itemchoice}
 	}
 \end{ex}
%========4
\begin{ex}%[0D2H1-3]
\immini{Cho bất phương trình $3x+y>3 \qquad(1)$.
	\choiceTF 
	{\True Bất phương trình $(1)$ là bất phương trình bậc nhất hai ẩn}
	{Cặp số $(x;y)=(1;0)$ là một nghiệm của bất phương trình $(1)$}
	{Bất phương trình $(1)$ có một nghiệm duy nhất}
	{\True Miền không được tô đậm (không kể bờ là đường thẳng $3x+y=3$) trong hình sau là miền nghiệm của bất phương trình $(1)$}
}{\begin{tikzpicture}[scale=0.9,line join=round, line cap=round,>=stealth,thick]
		\tikzset{every node/.style={scale=0.9}}
		\begin{scope}
			\clip (-1.5,-1) rectangle (2.5,4);
			\fill[pattern=north east lines] (-2,9)--(-2,-6)--(3,-6)--cycle;
			\draw (-0.33,4)--(1.33,-1)node [pos=0.45, above right] {$d$} ;
		\end{scope}
		\draw[->] (-1.5,0)--(2.5,0) node[below]{$x$};
		\draw[->] (0,-1)--(0,4) node[left]{$y$};
		\draw (0,0) node[below left]{$O$};
		\foreach \x in {1}
		\draw[thin] (\x,1pt)--(\x,-1pt) node [above right] {$\x$};
		\foreach \y in {3}
		\draw[thin] (1pt,\y)--(-1pt,\y) node [right] {$\y$};
	\end{tikzpicture}
} 	
 	\loigiai{
 		\begin{itemchoice}
 			\itemch Ta có $3x+y>3$ là bất phương trình bậc nhất hai ẩn dạng $ax+by>c$.
 			\itemch Thế cặp số $(x;y)=(1;0)$ vào bất phương trình $3x+y>3$ ta thấy $3 \cdot 1+ 0 = 3=3$ nên $(1;0)$ không thuộc miền nghiệm.
 			\itemch Vì bất phương trình bậc nhất hai ẩn luôn có một miền nghiệm thỏa mãn bất phương trình.
 			\itemch Đường thẳng $d \colon 3x+y=3$ đi qua hai điểm $(0;3)$ và $(1;0)$.\\
 			Lấy điểm $O(0;0) \notin d$. \\
 			Thay $x=0$ và $y=0$ vào bất phương trình $(1)$ ta được $0>3$, mệnh đề này sai.\\
 			Do đó miền nghiệm của bất phương trình $(1)$ là nửa mặt phẳng không kể bờ $d$, không chứa điểm $O(0;0)$.
 			\begin{center}
 				\begin{tikzpicture}[scale=1,line join=round, line cap=round,>=stealth,thick]
			\tikzset{every node/.style={scale=0.9}}
			\begin{scope}
				\clip (-2,-1) rectangle (3,4);
				\fill[pattern=north east lines] (-2,9)--(-2,-6)--(3,-6)--cycle;
				\draw (-0.33,4)--(1.33,-1)node [pos=0.45, above right] {$d$} ;
			\end{scope}
			\draw[->] (-2,0)--(3,0) node[below]{$x$};
			\draw[->] (0,-1)--(0,4) node[left]{$y$};
			\draw (0,0) node[below left]{$O$};
			\foreach \x in {1}
			\draw[thin] (\x,1pt)--(\x,-1pt) node [above right] {$\x$};
			\foreach \y in {3}
			\draw[thin] (1pt,\y)--(-1pt,\y) node [right] {$\y$};
		\end{tikzpicture}
 			\end{center}
 		\end{itemchoice}
 	}
 \end{ex}

%==============================================
\caukq
%========1
\begin{ex}%[0D2H1-2]
	\immini{Nửa mặt phẳng không bị gạch (không kể $d$) ở hình bên là miền nghiệm của bất phương trình có dạng $ax+by-4>0$. Giá trị của $a+b$ bằng bao nhiêu?}
	{
		\begin{tikzpicture}[scale=1,line join=round, line cap=round,>=stealth]
			\tikzset{every node/.style={scale=0.9}}
			\begin{scope}
				\clip (-1,-2.5) rectangle (3.5,1.5);
				\fill[pattern=north west lines] (-2,-4)--(-2,2.5)--(4.5,2.5)--cycle;
				\draw (3.5,1.5)--(-0.5,-2.5);
			\end{scope}
			\draw[->] (-1,0)--(3.5,0) node[below]{$x$};
			\draw[->] (0,-2.5)--(0,1.7) node[left]{$y$};
			\draw (0,0) node[below left]{$O$};
			\foreach \x in {2,3}
			\draw[thin] (\x,1pt)--(\x,-1pt) node [above] {$\x$};
			\foreach \y in {-2,-1}
			\draw[thin] (1pt,\y)--(-1pt,\y) node [left] {$\y$};
		\end{tikzpicture}
	}
	\shortans{$-2$}
	\loigiai{	
		Dựa vào hình ta thấy đường thẳng $d$ có hệ số góc khác $0$ nên gọi $(d)\colon y=ax+b$.\\
		Ta có $(2;0)\in (d)\Rightarrow 2a+b=0$ và $(0;-2)\in (d)\Rightarrow b=-2$.\\
		Vậy $(d)\colon y=x-2$ hay $x-y-2=0$.\\
		Mà $(0;0)$ không thuộc miền nghiệm $\Rightarrow x-y-2>0$ hay $2x-2y-4>0$.\\
		Do đó $a=2$; $b=2 \Rightarrow a+b=2+2=4$.
	}
\end{ex}
%========3
\begin{ex}%[0D2H1-3]
 	Bạn Nam tiết kiệm được $450$ nghìn đồng. Trong đợt ủng hộ các bạn học sinh đồng bào miền Trung bị lũ lụt vừa qua, bạn Nam đã ủng hộ $x$ tờ tiền loại $20$ nghìn đồng, $y$ tờ tiền loại $10$ nghìn đồng. Khi đó bất phương trình biểu diễn tổng số tiền mà bạn Nam đã ủng hộ có dạng $ax+y\le b$. Tính giá trị của biểu thức $P=2a-b$.
 	\shortans[oly]{$-41$}
 	\loigiai{
 		Tổng số tiền bạn Nam đã ủng hộ là $20 x+10 y$.\\
 		Vì bạn Nam chỉ có tất cả $450$ nghìn đồng nên tổng số tiền bạn Nam đã ủng hộ không thể vượt quá $450$ nghìn đồng. Vậy ta có bất phương trình $20 x+10 y \leq 450$ hay $2x+y \leq 45$.\\
 		Do đó $a=2$, $b=45$ nên $P=2\cdot 2-45=-41$.
 	}
 \end{ex}
%========4
\begin{ex}%[0D2N1-1]
Biết đường thẳng $d \colon 3x-4y-12=0$ đi qua hai điểm $A(0;a)$ và $B(b;0)$. Giá trị $a+b$ bằng bao nhiêu?
	\shortans[oly]{$1$}
	\loigiai{
	Thay tọa độ $(0;a)$ vào $d$ ta được $3 \cdot 0 - 4\cdot a-12=0 $ nên $a=-3$.\\
	Thay tọa độ $(b;0)$ vào $d$ ta được $3 \cdot b - 4\cdot 0-12=0 $ nên $b=4$.\\
	Vậy $a+b=-3+4=1$.
	}
\end{ex}
%========5
\begin{ex}%[0D2H1-3]
Bạn Lan mang $150\,000$ đồng đi nhà sách để mua một số quyển tập và bút. Biết rằng giá một quyển tập là $8\,000$ đồng và giá của một cây bút là $6\,000$ đồng. Bạn Lan có thể mua được tối đa bao nhiêu quyển tập nếu bạn đã mua $10$ cây bút.
	\shortans[oly]{$11$}
	\loigiai{
	Gọi $x$, $y$ lần lượt là số tập và số bút. Điều kiện $x>0$, $y>0$.
	\begin{itemize}
	\item Bất phương trình biểu diễn số tập và bút có thể mua được phụ thuộc vào số tiền mang theo là $8\,000x+6\,000y \leq 150\,000$.
	\item Bạn Lan có thể mua được tối đa số quyển tập nếu bạn đã mua $10$ cây bút là $8\,000x+6\,000 \cdot 10 \leq 150\,000 \Leftrightarrow x \leq 11{,}25$. 
	\item Vì $x$ nguyên dương nên số quyển tập tối đa bạn Lan mua được $11$ quyển.
	\end{itemize}
	}
\end{ex}
%========6
\cautl


\begin{bt}%[0D2N1-1]
	Cho bất phương trình hai ẩn $x-2y+6>0$.
	\begin{enumerate}
		\item[a)] $(0;0)$ có phải là một nghiệm của bất phương trình đã cho không?
		\item[b)] Biểu diễn miền nghiệm của bất phương trình đã cho trên mặt phẳng $Oxy$.
	\end{enumerate}	
	\loigiai{
		\begin{enumerate}
			\item[a)] Thế $(0;0)$ vào bất phương trình $x-2y+6>0$ ta được $0-2 \cdot 0+6=6>0$ (đúng). Do đó $(0;0)$ là một nghiệm của bất phương trình đã cho.
			\item[b)] Xét bất phương trình: $x-2y+6>0$.
			\immini{
				Vẽ đường thẳng $\Delta\colon x-2y+6=0$ đi qua hai điểm $A(0;3)$ và $B(-6;0)$.\\
				Xét gốc tọa độ $O(0;0)\notin \Delta$, ta có $0-2\cdot 0+6>0$.\\
				Do đó, miền nghiệm của bất phương trình là nửa mặt phẳng không kể bờ $\Delta$, chứa gốc tọa độ $O$ (miền không gạch chéo trên hình).}
			{\vspace{-0.5cm}
				\begin{tikzpicture}[line join=round, line cap=round,>=stealth,font=\footnotesize,scale=.7]
					\begin{scope}
						\clip (-7,-1) rectangle (1,4);
						\fill[pattern=north west lines] (-9,-1.5)--(-9,5)--(4,5)--cycle;
						\draw (1,3.5)--(-8,-1) node [pos=0.45, below] {$\Delta$};
					\end{scope}
					\draw[->] (-7,0)--(1,0) node[below]{$x$};
					\draw[->] (0,-1)--(0,4) node[right]{$y$};
					\draw (0,0) node[below left]{$O$};
					\foreach \x in {-6}
					\draw[thin] (\x,1pt)--(\x,-1pt) node [below] {$\x$};
					\foreach \y in {3}
					\draw[thin] (1pt,\y)--(-1pt,\y) node [below right] {$\y$};
			\end{tikzpicture}	}	
		\end{enumerate}
	}
\end{bt}
\begin{bt}%[0D2N1-2]
	Công ty viễn thông Viettel có gói cước Hi School tính phí là $1\,190$ đồng mỗi phút gọi nội mạng và $1\,390$ đồng mỗi phút gọi ngoại mạng. Một bạn học sinh đăng kí gói cước trên và sử dụng $x$ phút gọi nội mạng, $y$ phút gọi ngoại mạng trong một tháng. Viết bất phương trình bậc nhất hai ẩn $x$, $y$ để mô tả số tiền bạn đó phải trả trong một tháng ít hơn $100$ nghìn đồng.
	\loigiai{
		Trong một tháng, số tiền gọi nội mạng là $1\,190x$ đồng và số tiền gọi ngoại mạng là $1\,390y$ đồng.\\
		Tổng số tiền trong một tháng bạn học sinh phải trả là $1\,190x+1\,390y$.\\
		Để số tiền trong một tháng phải trả ít hơn $100$ nghìn đồng thì 
		\[
		1\,190x+1\,390y< 100\,000\Leftrightarrow 119y+139y< 10\,000.
		\]
	}
\end{bt}
%========BT9
\begin{bt}%[0D2N1-2]
	Bạn Cúc muốn pha hai loại nước cam. Để pha một lít nước cam loại $I$ cần $30$ g bột cam, còn một lít nước cam loại $II$ cần $20$ g bột cam. Gọi $x$ và $y$ lần lượt là số lít nước cam loại $I$ và $II$ pha chế được. Biết rằng Cúc chỉ có thể dùng không quá $100$ g bột cam. Hãy lập bất phương trình mô tả số lít nước cam loại $I$ và $II$ mà bạn Cúc có thể pha chế được và biểu diễn miền nghiệm của bất phương trình đó trên mặt phẳng toạ độ $Oxy$.
	\loigiai{
		\begin{itemize}
			\item Với $x$ lít $(x \geq 0)$ nước cam loại $I$ bạn Cúc cần sử dụng $20x$ g bột cam.
			\item Và $y$ lít $(y \geq 0)$ nước cam loại $II$ bạn Cúc cần sử dụng $30y$ g bột cam.
		\end{itemize}
		Suy ra tổng số bột cam bạn Cúc sử dụng để pha chế nước cam là: $20x+30y$.\\
		Ta có bất phương trình $20x+30y\le 100 \Leftrightarrow 2x+3y-10\le 0$.
		\immini{
			Vẽ đường thẳng $\Delta\colon 2x+3y-10=0$ đi qua hai điểm $A(2;2)$ và $B(5;0)$.\\
			Xét gốc tọa độ $O(0;0)\notin \Delta$, ta có $2\cdot 0+3\cdot 0-10<0$.\\
			Do đó, miền nghiệm của bất phương trình $2x+3y-10\le 0$ là nửa mặt phẳng kể cả bờ $\Delta$, chứa gốc tọa độ $O$ (miền không gạch chéo trên hình).}
		{\vspace{-0.5cm}
			\begin{tikzpicture}[line join=round, line cap=round,>=stealth,font=\footnotesize,scale=.7]
				\begin{scope}
					\clip (-1,-1) rectangle (6,4);
					\fill[pattern=north east lines] (-2,4.67)--(7,4.67)--(7,-1.33)--cycle;
					\draw (-1,4)--(6.5,-1) node [pos=0.55, below] {$\Delta$};
				\end{scope}
				\draw[->] (-1,0)--(6,0) node[below]{$x$};
				\draw[->] (0,-1)--(0,4) node[left]{$y$};
				\draw[dashed](2,0)|-(0,2);
				\draw (0,0) node[below left]{$O$};
				\foreach \x in {2,5}
				\draw[thin] (\x,1pt)--(\x,-1pt) node [below] {$\x$};
				\foreach \y in {2}
				\draw[thin] (1pt,\y)--(-1pt,\y) node [left] {$\y$};
		\end{tikzpicture}}
	}
\end{bt}
%========BT10
%==============================================
